\documentclass{article}
\usepackage{graphicx}
\usepackage{hyperref}
\usepackage{amsmath}
\usepackage{amssymb}
\usepackage{color,soul}
\usepackage{float}
\usepackage{csquotes}
\usepackage{latexsym}
\usepackage[left=1in, right=1in, top=1in, bottom=1in]{geometry}
\usepackage{pgfplots}
\usepackage{pgfplotstable}
\usepackage{listings}
\usepackage{xcolor}

\lstset{
    basicstyle=\small\ttfamily,
    frame=single,
    rulecolor=\color{gray!30},
    backgroundcolor=\color{gray!5}
}

\usepgfplotslibrary{units}
\pgfplotsset{
    compat=1.18,
    table/search path={../data},
    table/col sep=comma
}
\newcommand*{\code}[1]{\texttt{\detokenize{#1}}}
\begin{document}

\title{DATASCI 3ML3 Final Project\\{\Large{Wikipedia Article Classification}}}
\author{Liam Strilchuk\thanks{strilchl@mcmaster.ca}}
\date{April 20, 2026}
\maketitle

\tableofcontents
\cleardoublepage

\section{Overview and Motivation}

On Wikipedia, before articles are made visible to the public, they must go through a
review process. Reviewers, who are volunteers, must follow a 30+ step flowchart for each
published article;\footnote{See\url{https://upload.wikimedia.org/wikipedia/commons/f/f4/NPP_flowchart.svg}.}
this time-consuming process, along with a shortage of reviewers, means
there is often a months-long waiting time for articles to be reviewed (at the time of
writing, the number of articles waiting to be reviewed is over 17,000). Therefore, any
additional tools that could shorten the time taken to review an article would be incredibly
beneficial to the Wikipedia community.

For this project, we will create a model which can handle a step of this process. Articles
can be classified into one or more \textit{WikiProjects}; a WikiProject is a group of editors
focusing on a specific topic. For example, \href{https://en.wikipedia.org/wiki/Wikipedia:WikiProject_Canada}
{WikiProject Canada} is a group of (mainly Canadian) editors improving Canada-related articles.
Classifying articles into projects is extremely important, since it allows them to be targeted
for improvement by editors knowledgeable about the subject area; unclassified articles are
often neglected. However, since classifying is considered an ``optional'' step in the
reviewing process, many articles end up unclassified.

To this end, we will create a model to solve this multi-label classification problem. This
is different from multi-class classification, since an article can be part of zero, one, or
more WikiProjects. Since there are thousands of projects of varying size (for example,
\textit{Computing} has 36,000 articles; \textit{Artificial intelligence} has 1,600; and
\textit{Human-computer interaction} has only 184), we will focus on classifying only the
approximately 100 largest projects as a proof-of-concept. However, we will explore next
steps for a full model in the conclusion.

We will use Tensorflow/Keras to build and train the model, along with \code{scikit-learn}
for its natural language processing features (specifically, TF-IDF). In the remainder of this
report, we will discuss data acquisition and processing, creating a model, evaluation and
experimentation; finally, we will create an improved model based on these insights.

\section{Data Acquisition and Preprocessing}

To train our model, we need data in the format $(lead, projects)$, where the $lead$ is the
beginning article text, and $projects$ is a list containing the WikiProjects that article
has been assigned. While there is no existing dataset containing the exact data we need
for this model, Wikipedia is very open and provides a number of APIs which we can use to
collect the data ourselves. We will collect only the lead paragraphs of an article, instead
of its full content, as its projects can be determined from the first paragraph alone in
the vast majority of cases.

Wikipedia provides a directory of all WikiProjects, along with details about their
size.\footnote{See \url{https://en.wikipedia.org/wiki/Wikipedia:WikiProject_Directory/All}.}
To simplify the model, we will classify only the 25 WikiProjects with the largest number
of articles. Our first goal is to find a list of articles in each project; we can use
the \href{https://api.wp1.openzim.org}{WP 1.0 API} for this.

A major problem that we have to deal with here is the relative sizes of each category.
WikiProject Biography, for example, has nearly 2.2 million articles; while a smaller
WikiProject Lepidoptera has only around 100,000. This difference would be even more
pronounced if we were to classify more than 25 projects: the vast majority of projects
have fewer than 10,000 articles. If we were to use all possible data to train our model,
it would likely have a bias towards larger projects while potentially ignoring smaller
ones.\footnote{O. Olawale Awe, ``Computational Strategies for Handling Imbalanced Data in
Machine Learning''}

To combat this, we use a technique called undersampling. While there are fairly complex
algorithms that can be used for this purpose, we will use the simpler \textit{random
undersampling}, which will result in a more balanced dataset. The way this works in practice
is by training on only a portion of the available dataset for majority classes. So, we
will obtain only a random subset of articles in each project: approximately 2,000 from each.
This is likely far more than needed to train an effective model; in the Evaluation section
below, we will explore training models with varying amounts of available data.

After finding a list of the articles from each project, we will use Wikipedia's
\href{https://www.mediawiki.org/wiki/Extension:TextExtracts}{\code{extracts} API} to fetch
the lead section for each. Finally, since an article can be part of more than one project,
we go through each article and determine its other projects. Unfortunately, Wikipedia does
not provide a convenient API for this purpose; so instead, we use regular expressions to
parse the markup code for each article. We then save all of our collected data to
\code{all_data.csv}. An example row in this dataset is the following:

\begin{lstlisting}
Treaty of London (1867),"The Treaty of London (French: Traite de Londres), often
called the Second Treaty of London after the 1839 Treaty, granted Luxembourg full
independence...",WikiProject London|WikiProject Luxembourg|...
\end{lstlisting}

\section{Feature Engineering}

At the core of our model is the TF-IDF (term frequency-inverse document frequency) algorithm.
Unlike the simpler bag-of-words model, which only captures the frequency of a term over a
dataset, TF-IDF measures the relative importance of each term. We define $f_{t,d}$ as the number
of times term $t$ occurs in document $d$; we then define the \textit{term frequency} as

\begin{equation*}
	\text{TF}(t, d) = \frac{f_{t,d}}{\sum_{t'\in d}f_{t',d}}
\end{equation*}

Or, in simpler terms, TF represents the fraction of words that are $t$ in $d$. The other metric
we need to consider is \textit{inverse document frequency}. Let $D$ be the set of all documents
in our dataset:

\begin{equation*}
	\text{IDF}(t, D) = \log{\frac{|D|}{|\{d\in D : t\in d\}|}}
\end{equation*}

This measures the relative rarity of words across the dataset: higher values means the term
occurs in fewer documents. A logarithm is used because without it, a word that appears once
in a million documents would have an IDF score of 1,000,000, resulting in the algorithm weighing
it much higher than it should. We then multiply the TF and IDF values together to get the
final result:

\begin{equation*}
	\text{TF-IDF}(t, d, D) = \text{TF}(t, d)\times \text{IDF}(t, D)
\end{equation*}

For example, consider the word ``the''. As one of the most commonly used words in English, it
appears in approximately 96\% of the leads in our dataset; this would result in its IDF score
being $\log{\frac{100}{96}}\approx 0.02$, or an extremely low ``importance''. However, if we
consider the word ``Canadian'', which only appears in 4\% of leads, its IDF would be $\log
{\frac{100}{4}}\approx 1.4$.

For our model, we will use the \code{scikit-learn} library's \code{TfidfVectorizer} class.
This has a number of parameters to improve the algorithm's accuracy and performance:

\begin{itemize}
	\item \code{max_features} limits the number of terms, selecting only the top $N$ features
	by their term frequency. Setting this limit reduces the dimensionality of the data, which
	can reduce memory usage and prevent overfitting by excluding rare words.
	\item \code{stop_words} excludes common English words which do not provide a significant
	amount of information, helping the model focus on more important terms. This can also have
	the benefit of reducing the number of features and thus improving performance.
	\item \code{ngram_range} is a tuple $(min, max)$ that tells the vectorizer how many tokens
	in a row to consider as a single term. For example, with the input ``New York City'', the
	range $(1, 1)$ would produce \code{["New", "York", "City"]}, whereas $(1, 3)$ would also
	consider \code{"New York"} and \code{"New York City"}, allowing for more semantic context.
	However, this comes at the cost of a significantly higher number of features.
\end{itemize}

Importantly, when fitting the vectorizer to the data, we use only the training data, and not
any test data, to ensure there is no data leakage.

\section{Model Implementation}

\section{Evaluation and Insights}

\subsection{Experiment: Data Volume}

Since our model is classifying no more than the 100 largest WikiProjects, all of which have
at least 20,000 articles, our model is not suffering from a lack of training data. However,
should the model be expanded in the future to much smaller projects, some of which have only
hundreds of examples, this could become a problem. Therefore it would be beneficial to determine
how much data is needed to make accurate predictions.

Figure \ref{fig:metrics_vs_split} shows the results of the experiment. The chart shows the
classification metrics (Precision, Recall, and F1-Score) graphed against the portion of the
dataset used for training: between 1\% and 90\%. To ensure a fair comparison between these
radically different dataset sizes, instead of artificially constraining the number of epochs,
the Keras callback \code{EarlyStopping} was used to stop training when the validation accuracy
decreased. This ensures that models with less data do not overfit due to their limited sample size.

\begin{figure}[H]
    \centering
    \begin{tikzpicture}
        \begin{axis}[
            width=12cm,
            height=8cm,
            title={Performance Metrics vs. Data Split},
            xlabel={Split},
            ylabel={Score},
            xmin=0, xmax=0.9,
            ymin=0, ymax=1.0,
            xtick={0, 0.1, 0.2, 0.3, 0.4, 0.5, 0.6, 0.7, 0.8, 0.9},
            ytick={0, 0.1, 0.2, 0.3, 0.4, 0.5, 0.6, 0.7, 0.8, 0.9, 1.0},
            grid=major,
            grid style={dashed, gray!30},
            legend pos=outer north east,
            every axis plot/.append style={thick}
        ]
        \addplot[
            color=blue,
            mark=none
        ]
        table [x=split, y=precision, col sep=comma] {experiment_data_split.csv};
        \addlegendentry{Precision}
        \addplot[
            color=red,
            mark=none
        ]
        table [x=split, y=recall, col sep=comma] {experiment_data_split.csv};
        \addlegendentry{Recall}
        \addplot[
            color=green!70!black,
            mark=none
        ]
        table [x=split, y=f1-score, col sep=comma] {experiment_data_split.csv};
        \addlegendentry{F1-Score}
        \end{axis}
    \end{tikzpicture}
    \caption{Precision, Recall, and F1-score evaluated across different data splits.}
    \label{fig:metrics_vs_split}
\end{figure}

While the model's performance when using less than 10\% of the dataset is extremely poor,
the accuracy begins to level out around 20\%, and by 50\% we experience diminishing returns
and any improvement becomes very marginal. While increasing the amount of data used does
result in an improvement, even if small, it comes at a cost: the training time increases
significantly when more data is used. For example, when all available data was used to train
a model, the training time was roughly 3x that of a model using only 30\%; however, the
improvement in accuracy was only 2\%. Thus it may not be worth it to always use the
largest possible dataset.

\subsection{Experiment: Dropout}

Dropout is a regularization technique used in neural networks to prevent overfitting. It
entails deactivating a random portion of neurons when training, forcing the network to learn
robust patterns instead of memorizing the training data. In this experiment, we train models
with different levels of dropout and examine their training and validation loss throughout
the process. This experiment uses models that have low, medium, and high levels of dropout
(0.0, 0.3, and 0.6 respectively).

\begin{figure}[H]
    \centering
    \begin{tikzpicture}
        \begin{axis}[
            width=12cm,
            height=8cm,
            title={Model Loss (Training vs. Validation)},
            xlabel={Epoch},
            ylabel={Loss},
            xmin=1, xmax=15,
            ymin=0, ymax=0.2,
            xtick={2, 4, 6, 8, 10, 12, 14},
            grid=major,
            grid style={dashed, gray!30},
            legend pos=outer north east,
            every axis plot/.append style={thick}
        ]
        \addplot[
            color=blue,
            mark=none
        ]
        table [x=epoch, y=low_loss, col sep=comma] {experiment_dropout.csv};
        \addlegendentry{Loss (low)}
        \addplot[
            color=red,
            mark=none
        ]
        table [x=epoch, y=low_val_loss, col sep=comma] {experiment_dropout.csv};
        \addlegendentry{Val Loss (low)}
        \addplot[
            color=yellow!80!black, 
            mark=none
        ]
        table [x=epoch, y=med_loss, col sep=comma] {experiment_dropout.csv};
        \addlegendentry{Loss (med)}
        \addplot[
            color=green!70!black,
            mark=none
        ]
        table [x=epoch, y=med_val_loss, col sep=comma] {experiment_dropout.csv};
        \addlegendentry{Val Loss (med)}
        \addplot[
            color=orange,
            mark=none
        ]
        table [x=epoch, y=high_loss, col sep=comma] {experiment_dropout.csv};
        \addlegendentry{Loss (high)}
        \addplot[
            color=cyan!80!black,
            mark=none
        ]
        table [x=epoch, y=high_val_loss, col sep=comma] {experiment_dropout.csv};
        \addlegendentry{Val Loss (high)}
        \end{axis}
    \end{tikzpicture}
    \caption{Training and validation loss across epochs for different dropout rates.}
    \label{fig:dropout_loss}
\end{figure}

Figure \ref{fig:dropout_loss} shows the results of the experiment. As expected, the model
with no dropout has its training loss quickly drop to near zero while its validation loss
steadily rises, indicating overfitting. While the medium and high dropout rates smooth
this curve and delay the divergence, they fail to prevent the validation loss from
increasing. All models achieve their absolute lowest validation loss very early, around
epoch 2-3, showing the importance of the early stopping technique. Since even a high level
of dropout fails to prevent overfitting, this suggests that the model possesses excess
capacity, and reducing the size of the network's layers may be more effective at preventing
overfitting.

\subsection{Experiment: Max Features}

The \code{max_features} parameter of the \code{TfidfVectorizer} class was discussed in
Section 3. Optimizing its value can significantly impact both the training time and
performance of a model. Too low, and the model does not have enough data to accurately
classify text; too high, and the memory usage and training time increase, and we risk
providing the model with thousands of ``filler'' words which do not matter for classification.

This experiment tests the performance metrics of models based on the maximum number
of features. Figure \ref{fig:metrics_vs_max_features} below shows the results.

\begin{figure}[H]
    \centering
    \begin{tikzpicture}
        \begin{semilogxaxis}[
            width=12cm,
            height=8cm,
            title={Performance Metrics vs. Max Features},
            xlabel={Max Features},
            ylabel={Score},
            ymin=0.4, ymax=1.0,
            xmin=50, xmax=30000,
            xtick={50, 100, 500, 1000, 5000, 10000}, 
            log ticks with fixed point,
            grid=major,
            grid style={dashed, gray!30},
            legend pos=outer north east,
            every axis plot/.append style={thick}
        ]
        \addplot[
            color=blue,
            mark=none
        ]
        table [x=count, y=precision, col sep=comma] {experiment_max_features.csv};
        \addlegendentry{Precision}
        \addplot[
            color=red,
            mark=none
        ]
        table [x=count, y=recall, col sep=comma] {experiment_max_features.csv};
        \addlegendentry{Recall}
        \addplot[
            color=green!70!black,
            mark=none
        ]
        table [x=count, y=f1-score, col sep=comma] {experiment_max_features.csv};
        \addlegendentry{F1-Score}
        \end{semilogxaxis}
    \end{tikzpicture}
    \caption{Performance metrics evaluated across a count of max features (logarithmic scale).}
    \label{fig:metrics_vs_max_features}
\end{figure}

Even 100 features is enough to achieve a reasonable degree of accuracy. The increase in model
performance becomes marginal once the number of features reaches 1,000; after 10,000, there is
no improvement whatsoever. This is likely because Wikipedia articles, despite covering a vast
number of topics, can be categorized using a relatively small vocabulary: the words ``album''
or ``Canadian'' can tell us, for example, that an article should be put into the ``Albums''
or ``Canada'' WikiProjects respectively, without a large amount of other context needed.

\subsection{Experiment: Model Size}

Finally, we will run an experiment comparing models with different sizes of hidden layers.
We measure the layer size as a factor of our arbitrarily picked initial values of 512 and 256
for the first and second layers respectively. The more neurons a network has, the more capacity
it has for learning complex patterns; however, this comes at the cost of increased memory
usage, training time, and potential for overfitting. For example, a network with dense layers
of sizes 1,024 and 512 has 524,288 trainable parameters between them; if the sizes are reduced
by a factor of 8 to 128 and 64, this lowers to only 8,192.

\begin{figure}[H]
    \centering
    \begin{tikzpicture}
        \begin{axis}[
            width=12cm,
            height=8cm,
            title={Performance Metrics vs. Size Factor},
            xlabel={Size Factor},
            ylabel={Score},
            xmin=0, xmax=1.4,
            ymin=0, ymax=1.0,
            xtick={0.25, 0.50, 0.75, 1.00, 1.25},
            ytick={0, 0.25, 0.50, 0.75, 1.00},
            grid=major,
            grid style={dashed, gray!30},
            legend pos=outer north east,
            every axis plot/.append style={thick}
        ]
        \addplot[
            color=blue,
            mark=none
        ]
        table [x=factor, y=precision, col sep=comma] {experiment_model_size.csv};
        \addlegendentry{precision}
        \addplot[
            color=red,
            mark=none
        ]
        table [x=factor, y=recall, col sep=comma] {experiment_model_size.csv};
        \addlegendentry{recall}
        \addplot[
            color=green!70!black,
            mark=none
        ]
        table [x=factor, y=f1-score, col sep=comma] {experiment_model_size.csv};
        \addlegendentry{f1-score}
        \end{axis}
    \end{tikzpicture}
    \caption{Performance metrics evaluated across different model size factors.}
    \label{fig:metrics_vs_factor}
\end{figure}

Figure \ref{fig:metrics_vs_factor} shows the results of this experiment. We can see that
our initial estimation of 512/256 was significantly more than needed. The performance metrics
increase sharply until a factor of 0.15, at which point they even out and remain constant
thereafter. Scaling the model any more than 0.25 is waste of resources, as it does not
provide any performance gains.

\subsection{Final Model}

We will now create an optimized model based on the results of the previous experiments.
Firstly, our testing with data volume revealed that there is a consistent (although gradual)
increase in performance

\section{Conclusion}

\subsection{Next Steps}

\begin{itemize}
	\item Partial random under-sampling\footnote
	{Nicolás García-Pedrajas, ``Partial random under/oversampling for multilabel problems''}
	\item Pretrained models or word embeddings
	\item Categories instead of projects
\end{itemize}

\section{Code Organization}

The code for this project is split between Python modules and Jupyter notebooks: testing,
experiment, and data collection code is in notebooks, while the modularized \code{Model}
class is in a Python file. The description for each file is as follows:

\begin{itemize}
	\item \code{data_collection.ipynb} uses the steps described in Section 2 to collect data:
	it uses Wikipedia's API to collect articles in each project, find the lead text for each,
	find other WikiProjects for each article, and then saves the result to \code{all_data.csv}.
	\item \code{create_model.ipynb} walks through the steps for creating a model: using a
	binarizer to turn WikiProjects into vectors, splitting test and training data, vectorizing
	the text, and creating, training, and saving the model.
	\item \code{model.py} generalizes the code from this notebook to make a modular interface
	with a \code{Model} class. This allows us to pass in additional parameters to run
	experiments: dataset size, training duration, dropout, etc. It also provides \code{save},
	\code{load}, and \code{get_report} functions, so we can easily tell how a model is
	performing, or load a previously trained model.
	\item \code{evaluation.ipynb} contains a function, \code{run_eval_thread}, which uses the
	\text{Model} class to create and train a model with the given parameters. It does this on
	a separate thread, to ensure memory is reclaimed once it is done. It then runs the model's
	\code{get_report} function, and returns statistics.
\end{itemize}

\section{References}

\end{document}